\documentclass[11pt]{article}
\newcommand{\paso}{4}
\usepackage[T1]{fontenc}
\usepackage[utf8]{inputenc}
\usepackage[spanish,es-nodecimaldot,es-noquoting,es-lcroman]{babel}
\spanishplainpercent
\usepackage{lmodern}
\usepackage[a4paper,margin=21mm,headheight=15pt]{geometry}
\usepackage{amsmath,amssymb,bm}
\usepackage{booktabs,array}
\usepackage{tikz}
\usetikzlibrary{arrows.meta,positioning,calc}
\usepackage[most]{tcolorbox}
\usepackage{microtype}
\usepackage{fancyhdr}
\usepackage{hyperref}
\hypersetup{colorlinks=true,linkcolor=blue!45!black,urlcolor=blue!45!black}
\setlength{\parindent}{0pt}
\setlength{\parskip}{5pt}
\setlength{\emergencystretch}{2em}
\pagestyle{fancy}
\fancyhf{}
\fancyhead[L]{\small BEM del cilindro: paso \paso}
\fancyhead[R]{\small Guía matemática independiente}
\fancyfoot[C]{\thepage}
\newcommand{\vct}[1]{\bm{#1}}
\newcommand{\dd}{\,\mathrm{d}}
\newcommand{\ii}{\mathrm{i}}
\newcommand{\norm}[1]{\left\lVert#1\right\rVert}
\newcommand{\VP}{\operatorname{VP}}
\newcommand{\tr}{\operatorname{tr}}
\newtcolorbox{idea}[1]{breakable,colback=blue!3,colframe=blue!50!black,
  title={#1},fonttitle=\bfseries,boxrule=0.7pt}
\newtcolorbox{alerta}[1]{breakable,colback=orange!4,colframe=orange!65!black,
  title={#1},fonttitle=\bfseries,boxrule=0.7pt}
\newtcolorbox{resultado}[1]{breakable,colback=teal!3,colframe=teal!55!black,
  title={#1},fonttitle=\bfseries,boxrule=0.7pt}
\newcommand{\inicio}[2]{%
 {\small\color{blue!55!black}\bfseries SERIE: LOS SIETE PASOS MATEMÁTICOS DEL BEM}\par
 {\LARGE\bfseries Paso \paso. #1\par}\vspace{3mm}
 \begin{idea}{¿Qué busca este paso?}#2\end{idea}}
\newcommand{\mapa}[3]{%
 \begin{center}
 \begin{tikzpicture}[node distance=4mm,
   caja/.style={draw=blue!50!black,fill=blue!3,rounded corners=2pt,
     text width=13cm,align=center,inner sep=6pt},
   >={Stealth[length=2mm]}]
 \node[caja] (a){#1};
 \node[caja,below=of a] (b){#2};
 \node[caja,below=of b] (c){#3};
 \draw[->,thick,blue!60!black] (a)--(b);
 \draw[->,thick,blue!60!black] (b)--(c);
 \end{tikzpicture}
 \end{center}}
\newcommand{\base}{%
 \smallskip{\footnotesize\color{black!65}
 Modelo de referencia: cilindro sólido, material homogéneo e isótropo,
 deformación plana, pequeñas deformaciones, sin fuerzas de volumen;
 amplitudes complejas con \(e^{+\ii\omega t}\), \(\omega>0\).}\par}
\newcommand{\fuentes}{%
 \par\medskip\begingroup\footnotesize
 \textbf{Referencia y continuación.}
 Este documento desarrolla una etapa del esquema
 \texttt{esquema\_ecuaciones\_bem.tex}. Las ecuaciones y los ejemplos se
 contrastaron con \texttt{bem\_cilindro.py}, ambos en la carpeta superior.
 La guía general y sus deducciones ampliadas se encuentran en
 \texttt{../guia\_bem/guia\_bem.pdf}. No se modifica la implementación.
 \par Para recompilar este documento, mantener \texttt{estilo.tex} en
 la misma carpeta y ejecutar \texttt{pdflatex} sobre su archivo
 \texttt{.tex}, desde esa carpeta.\par\endgroup}

\begin{document}
\inicio{Los núcleos dinámicos}{Construir las funciones de influencia
\(U\) y \(T\) que hacen utilizable la identidad de frontera. Para cualquier
par fuente--integración, deben proporcionar la respuesta auxiliar del
material y su tracción, con frecuencia, normal y convención temporal
coherentes.}
\base
\mapa{Material: \(\lambda^*,\mu^*,\rho\);\quad frecuencia \(\omega\)}
{Distancia y orientación: \(r=\norm{\vct x-\vct\xi}\), \(\vct e=(\vct x-\vct\xi)/r\), \(\vct n(\vct x)\)}
{Dos matrices \(2\times2\): \(U(\vct x,\vct\xi;\omega)\) y \(T(\vct x,\vct\xi;\omega)\)}

\section*{1. La pregunta auxiliar}
Se aplica una fuerza puntual unitaria por unidad de longitud axial en
\(\vct\xi\), en un medio infinito del mismo material. Para la fuerza en
dirección \(a\), el campo auxiliar es
\[
 v_j(\vct x)=U_{aj}(\vct x,\vct\xi).
\]
No es la presión real del cilindro: es un experimento matemático con el que
se construye una base de influencias.

\begin{idea}{La convención de índices que se mantendrá}
El primer índice \(a\) identifica la fuerza auxiliar o ecuación de colocación;
el segundo \(j\), la componente del dato integrado.
\[
 u_a(\vct\xi)=\int_\Gamma U_{aj}t_j\,\dd\Gamma
              -\int_\Gamma T_{aj}u_j\,\dd\Gamma.
\]
La tracción \(T_{aj}\) pertenece a la \textbf{fila} \(U_{a\cdot}\).
No cambiarla por la tracción de una columna sin trasponer la convención.
\end{idea}

\section*{2. Distancia local frente al radio del cilindro}
\[
 r=\norm{\vct x-\vct\xi},\qquad e_i=\frac{x_i-\xi_i}{r},\qquad
 s=\vct e\cdot\vct n.
\]
\(r\) cambia para cada par de puntos y puede tender a cero.
El radio \(R=6\) m es una propiedad fija del dominio.
\(s\) es un coseno orientado: no es otra distancia ni un jacobiano.

\newpage
\section*{3. Por qué aparecen dos funciones de onda}
El operador elástico permite separar contribuciones longitudinales y
transversales. Sus números de onda son
\[
 k_s=\omega/c_s,\qquad k_p=\omega/c_p.
\]
Una función fundamental escalar adecuada a \(e^{+\ii\omega t}\) es
\[
 g_\alpha(r)=\frac{1}{2\pi}K_0(\ii k_\alpha r),\qquad
 (\nabla^2+k_\alpha^2)g_\alpha=-\delta(\vct x-\vct\xi),
 \quad\alpha=s,p.
\]
\(K_n\) denota Bessel modificada de segunda especie. Su argumento es
complejo. Por eso su presencia no significa que se esté resolviendo un
problema estático de difusión.

\section*{4. De funciones escalares al tensor elástico}
La combinación que satisface la ecuación elástica con fuente unitaria es
\[
 \boxed{U_{aj}=\frac1{\mu^*}
 \left[\delta_{aj}g_s+\frac1{k_s^2}
 \partial_a\partial_j(g_s-g_p)\right].}
\]
El primer término aporta la parte transversal, pero por sí solo no
satisface todo el operador elástico. El término de derivadas corrige su
contribución longitudinal.

Para una función radial \(f(r)\):
\[
 \partial_i f=f'e_i,\qquad
 \partial_i\partial_jf=
 \left(f''-\frac{f'}r\right)e_ie_j+\frac{f'}r\delta_{ij}.
\]
Así, un tensor construido a partir de una función radial tiene dos piezas:
una proporcional a la identidad y otra a \(\vct e\vct e^{\mathsf T}\).

\begin{resultado}{La forma que conviene para evaluar el núcleo}
Definiendo
\[
 z_s=\ii\omega r/c_s,\qquad z_p=\ii\omega r/c_p,\qquad
 \gamma=c_s/c_p,\qquad z_p=\gamma z_s,
\]
las derivadas de Bessel se agrupan en dos funciones escalares \(\psi,\chi\).
Los argumentos \(z_s,z_p\) son adimensionales.
\end{resultado}

\newpage
\section*{5. La fórmula explícita de desplazamiento}
Con \(K_0'(z)=-K_1(z)\) y
\(K_2(z)=K_0(z)+2K_1(z)/z\), al sustituir el hessiano radial se obtiene
\[
 \psi(r)=K_0(z_s)+\frac{K_1(z_s)-\gamma K_1(z_p)}{z_s},
 \qquad
 \chi(r)=K_2(z_s)-\gamma^2K_2(z_p),
\]
\[
 \boxed{U_{aj}=\frac{\psi\,\delta_{aj}-\chi\,e_ae_j}{2\pi\mu^*}.}
\]
La matriz diádica que aparece aquí es
\[
 \vct e\vct e^{\mathsf T}
 =\begin{pmatrix}e_x^2&e_xe_y\\e_xe_y&e_y^2\end{pmatrix}.
\]
De ahí se leen las cuatro componentes:
\[
 U=\frac1{2\pi\mu^*}
 \begin{pmatrix}
 \psi-\chi e_x^2&-\chi e_xe_y\\
 -\chi e_xe_y&\psi-\chi e_y^2
 \end{pmatrix}.
\]

\section*{6. Interpretar la orientación sin perderse en índices}
Si fuente y punto de integración están alineados horizontalmente,
\(\vct e=(1,0)\), entonces
\[
 U=\frac1{2\pi\mu^*}
 \begin{pmatrix}\psi-\chi&0\\0&\psi\end{pmatrix}.
\]
En esa base, \(\psi-\chi\) controla la componente paralela a la separación
y \(\psi\), la perpendicular. Para una separación oblicua aparecen
componentes cruzadas.

\begin{alerta}{Simetría de \(U\) no implica simetría de todos los operadores}
Este \(U\) es simétrico respecto a sus dos componentes. Sin embargo,
el operador de tracción incluye una normal y derivadas con una orientación
definida; en general \(T_{aj}\ne T_{ja}\).
Tampoco se deduce que las matrices globales de colocación sean simétricas.
\end{alerta}

\section*{7. Unidades}
\(\psi,\chi\) son adimensionales y \([U]=\mathrm{Pa}^{-1}\).
En la identidad BEM,
\([U\,\vct t\,\dd\Gamma]=\mathrm{Pa}^{-1}\mathrm{Pa}\,\mathrm m=\mathrm m\),
como corresponde a un desplazamiento.

\newpage
\section*{8. De \(U\) a su gradiente}
La tracción requiere esfuerzos, los esfuerzos requieren deformaciones y
estas requieren derivadas espaciales. La cadena es
\[
 U\ \longrightarrow\ \partial_kU\
 \longrightarrow\ \varepsilon^{\mathrm{aux}}\
 \longrightarrow\ \sigma^{\mathrm{aux}}\
 \longrightarrow\ T=\sigma^{\mathrm{aux}}\vct n.
\]
Las identidades geométricas son
\[
 \partial_kr=e_k,\qquad
 \partial_ke_j=\frac{\delta_{jk}-e_je_k}{r}.
\]
Al derivar \(U\):
\[
\begin{aligned}
 \partial_kU_{aj}=\frac1{2\pi\mu^*}\bigg[
 &\psi'e_k\delta_{aj}-\chi'e_ke_ae_j\\
 &-\frac{\chi}{r}
 (\delta_{ak}e_j+\delta_{jk}e_a-2e_ae_je_k)\bigg].
\end{aligned}
\]
Las primas indican derivadas con respecto a \(r\), no con respecto al
argumento de Bessel ni al parámetro local del elemento.

\section*{9. La regla de la cadena en las funciones radiales}
Sean \(a_s=\ii\omega/c_s\), \(a_p=\ii\omega/c_p\). Entonces
\(z_s=a_sr\), \(z_p=a_pr\), y
\[
\begin{aligned}
 \psi'={}&a_sK_0'(z_s)
 +\frac{a_sK_1'(z_s)-\gamma a_pK_1'(z_p)}{z_s}\\
 &-\frac{a_s[K_1(z_s)-\gamma K_1(z_p)]}{z_s^2},\\
 \chi'={}&a_sK_2'(z_s)-\gamma^2a_pK_2'(z_p).
\end{aligned}
\]
En estas expresiones, \(K_n'(z)\) sí significa derivada respecto a \(z\).
Omitir \(a_s\) o \(a_p\) cambia tanto las unidades como la frecuencia
del núcleo de tracción.

\begin{idea}{Una comprobación local antes del ensamblaje}
Para \(r\) no demasiado pequeño, se puede comparar la derivada analítica
con diferencias finitas de \(U\) y luego aplicar Hooke. Una buena solución
del sistema lineal no detecta por sí sola un error en estas derivadas.
\end{idea}

\newpage
\section*{10. Aplicar Hooke y proyectar con la normal}
Para la fila auxiliar \(U_{a\cdot}\), la componente de tracción es
\[
 \boxed{T_{aj}=\lambda^*n_j\partial_kU_{ak}
       +\mu^*n_k(\partial_kU_{aj}+\partial_jU_{ak}).}
\]
Esta es la forma inequívoca de fijar la convención de índices.
Agrupando las derivadas y usando \(s=e_kn_k\):
\[
\boxed{\begin{aligned}
 T_{aj}=\frac1{2\pi}\bigg[
 &\left(\psi'-\frac{\chi}{r}\right)(\delta_{aj}s+e_jn_a)\\
 &-\frac{2\chi}{r}(e_an_j-2e_ae_js)
 -2\chi'e_ae_js\\
 &+\frac{\lambda^*}{\mu^*}
   \left(\psi'-\chi'-\frac{\chi}{r}\right)e_an_j
 \bigg].
\end{aligned}}
\]
La normal es la del punto de integración sobre el elemento.
El núcleo depende de ella linealmente; invertir la normal cambia su signo.
\([T]=\mathrm m^{-1}\), así que
\([T\,\vct u\,\dd\Gamma]=\mathrm m\).

\section*{11. Separar la parte estática singular}
La tracción estática de Kelvin, para la misma convención, es
\[
\boxed{\begin{aligned}
 T^s_{aj}=-\frac1{4\pi(1-\nu)r}
 \bigl[&s\big((1-2\nu)\delta_{aj}+2e_ae_j\big)\\
       &-(1-2\nu)(e_an_j-e_jn_a)\bigr].
\end{aligned}}
\]
Su dependencia en \(\mu^*\) se cancela: para este material depende de
\(\nu\), la geometría y la distancia, pero no de la frecuencia.

\begin{resultado}{La identidad que permite regularizar}
\[
 T=T^s+(T-T^s).
\]
No es una aproximación física: se suma y resta el mismo núcleo.
La aproximación aparece después al integrar numéricamente. La parte
estática diagonal se construye con la identidad rígida del paso 6.
\end{resultado}

\newpage
\section*{12. Qué sucede al acercarse a la fuente}
Para frecuencia fija y \(r\to0\), las expansiones de Bessel dan
\[
 \chi(r)=\frac{\gamma^2-1}{2}+O(r^2\log r),\qquad
 \psi(r)=-\frac{1+\gamma^2}{2}\log(r/r_0)+C(\omega)
          +O(r^2\log r).
\]
\(r_0\) es una longitud de referencia; un cambio de \(r_0\) modifica la
constante, no el orden singular. Por tanto,
\[
 U=O(\log r),\qquad T=O(1/r),\qquad T-T^s=O(r\log r).
\]
\begin{itemize}
 \item La singularidad logarítmica es integrable en un elemento de borde.
 \item La parte \(1/r\) requiere el tratamiento de valor principal.
 \item La corrección dinámica pierde esa singularidad fuerte; no por ello
 es numéricamente exacta ni necesariamente analítica en la fuente.
\end{itemize}

\section*{13. Un límite que permite comprobar la evaluación}
Con \(\nu=0.25\), \(\gamma^2=1/3\), de modo que
\[
 \chi\longrightarrow-\frac13,\qquad \chi'\longrightarrow0.
\]
A \(\omega=40\) rad/s, para el material de referencia y \(r=0.1\) m:
\[
 \psi\simeq2.148475-\ii1.013374,\qquad
 \chi\simeq-0.334014+\ii0.000338.
\]
Estos valores se obtienen directamente de las funciones radiales del modelo.
No se evalúa la fórmula en \(r=0\).

\begin{alerta}{Cancelación numérica cerca de la fuente}
En \(\chi=K_2(z_s)-\gamma^2K_2(z_p)\) se restan cantidades grandes para
obtener una cantidad finita. La aritmética de precisión finita puede perder
dígitos, y las derivadas pueden ser todavía más sensibles. Más puntos de
Gauss acercan algunas evaluaciones a la fuente y no garantizan mejora.
Las expansiones asintóticas ofrecen una vía de estabilización, pero no se
han añadido ni supuesto en la implementación existente.
\end{alerta}
\section*{14. Qué sale de este paso}
Para cada fuente y punto de cuadratura se conocen \(U,T,T^s\).
El paso 5 indica dónde evaluarlos y con qué pesos; el paso 6 suma sus
contribuciones para construir las matrices del problema real.
\fuentes
\end{document}
